\section{Hamiltonians with quantitative unitary dynamics}
\label{sec:dynamics}

We use units in which $\hbar=1$. An unbounded observable is specified by its
\emph{graph}: a state vector $x$ in its domain comes with a second Hilbert
name $y$ and witnesses of the equations that say $y=Hx$. The name of $y$
supplies information that can be stronger than a bound on $\|y\|$.
For example, differentiating a trajectory in Hilbert norm requires the
tail of its derivative to be small at a known cutoff.

A strong derivative certificate for a curve $X$ at $t$ consists of a
Hilbert name $z$ and a PR function $d(k)$ such that
\[
 \|X(t+h)-X(t)-hz\|\le 2^{-k}|h|
 \quad\text{whenever } |h|\le 2^{-d(k)} .
\]
One may test rational increments and extend to named increments by
the continuity estimates, using rational approximants inside the closed
interval. To compare two proposed derivatives, substitute the positive
rational $h=2^{-\max(d_1(k),d_2(k))}$; rational cancellation gives
$\|z_1-z_2\|\le2^{1-k}$ for arbitrary $k$.

\subsection{Energy coordinates}

Let $e_j$ be the canonical basis of $\ell^2(\mathbb N)$. Supply real
energy names $\lambda_j$ uniformly in $j$, with a PR upper bound
$\Lambda(N)\ge 1+\max_{j<N}|\lambda_j|$; an integer bound of this kind
can always be computed from the stage-zero approximations. Define
\[
 (x,y)\in D \ \Longleftrightarrow\ (\forall j)\ y_j=\lambda_jx_j,
 \qquad (U(t)x)_j=e^{-it\lambda_j}x_j .
\]
Here and below a universal statement about the coordinates abbreviates
tests at arbitrary finite index and precision. Write $P_N$ for the first
$N$ coordinate projection and $Q_N=I-P_N$. A finite approximation $v$
to $x$ within $\epsilon$, supported below $N$, proves
$\|Q_Nx\|\le\epsilon$ by contraction.

\begin{theorem}[Diagonal graph and generator]\label{thm:diag-energy}
$D$ is a densely defined, closed, self-adjoint graph. The maps $U(t)$
form a strongly continuous unitary group whose strong derivative at zero
has graph $-iD$. For every graph pair $(x,y)$, $U(t)x$ has derivative
$-iU(t)y$ at $t$. All constructions have PR realizers.
\end{theorem}
\begin{proof}
Uniqueness and closedness follow by applying a coordinate functional to
the approximants: equal coordinates give equal vectors by truncating both
tails. Every basis vector is in the domain, and the graph truncations
$(P_Nx,P_Ny)$ converge in both components with an explicit joint cutoff.
The adjoint graph consists of pairs $(x,z)$ with
$\langle Dv,x\rangle=\langle v,z\rangle$ on graph inputs $v$
(inner products are linear in the second slot).
Testing $v=e_j$ forces $z_j=\lambda_jx_j$.
Conversely the equalities on finite sums pass to graph inputs by the joint
cutoff and Cauchy--Schwarz. The resolvents $(D\pm i)^{-1}$ are diagonal
multipliers by $(\lambda_j\pm i)^{-1}$. Both these coefficients and
$\lambda_j/(\lambda_j\pm i)$ have modulus at most one. This supplies
names for the solution and its graph partner for every right-hand side.

For computation of a phase or a finite matrix exponential, the elementary
bound
\begin{equation}\label{eq:exp-pr}
 \sum_{r=m}^\infty {a^r\over r!}
 \le 2^{-m}\sum_{r=0}^\infty{(2a)^r\over r!}
 \le 2^{4L-m}\qquad(0\le a\le L,\ L\in\mathbb N)
\end{equation}
follows by multiplying the series with positive coefficients and using
$\sum 1/r!\le3$, hence $\exp(2L)\le3^{2L}$. These series identities can
equally be verified by the binomial convolution of finite coefficients
and the geometric tail once successive terms have ratio at most $1/2$.
Absolute product and adjoint identities give
$e^{is}e^{it}=e^{i(s+t)}$ and $|e^{it}|=1$ for real $s,t$.
Integrating the convergent series on a bounded interval gives
\[
 |e^{is}-1|\le |s|,\qquad |e^{is}-1-is|\le s^2/2 .
\]
For error $\epsilon=2^{-k}$, take a finite $v$ within $\epsilon/4$ of
$x$. On its support use uniform phase errors at most
$\epsilon/(4(1+\|v\|))$ and form a finite output to accuracy
$\epsilon/4$. The parameter $L$ in \eqref{eq:exp-pr} is bounded from
$|t|$ and $\Lambda(N)$. For
$|t|\le T$ and $|\lambda|\le M$, rational approximations to $t,\lambda$
within $\rho\le1$ give a product error at most $(T+M+1)\rho$;
the unit-phase map is $1$-Lipschitz. All norms appearing in precision choices may be
rounded upwards using rational upper estimates. Coordinate multiplication
proves the group and adjoint laws. On a head and tail one obtains
\[
 \|(U(h)-I)x\|\le |h|\Lambda(N)\|x\|+2\|Q_Nx\|.
\]
This is a strong continuity modulus at every $t$. For a graph pair the
first-order remainder satisfies
\[
 \|U(h)x-x+ihy\|
 \le |h|\bigl(2\|Q_Ny\|+|h|\Lambda(N)^2\|x\|/2\bigr).
\]
Indeed, on each tail coordinate the first-order difference is bounded
by $2|h||y_j|$, and on the head the second scalar bound applies.
Choose $N$ from the name of $y$ and then a sufficiently small dyadic
bound on $|h|$. Conversely, a named derivative $z$ implies, by the same
scalar expansion at each fixed $j$, that $z_j=-i\lambda_jx_j$.
Thus $(x,iz)$ is a graph pair. Translating time gives the remaining
claim.
\end{proof}

\subsection{A number-basis oscillator}

On finite sequences put $ae_0=0$,
$ae_n=\sqrt n\,e_{n-1}$ for $n>0$, and
$a^\dagger e_n=\sqrt{n+1}\,e_{n+1}$. The square roots are obtained
by bounded bisection. Set
\[
 Q=(a+a^\dagger)/\sqrt2,\qquad
 P=(a-a^\dagger)/(i\sqrt2),\qquad
 H=(P^2+Q^2)/2 .
\]
The core of finite sequences is invariant under every displayed
operator. Inner products of basis vectors give the stated adjunction.
Acting on $e_n$ gives $aa^\dagger e_n=(n+1)e_n$ and
$a^\dagger a e_n=ne_n$, including $n=0$. Expanding products therefore
proves on the core
\[
 [Q,P]=iI,\qquad H=a^\dagger a+\tfrac12 I,\qquad
 U(t)e_n=e^{-it(n+1/2)}e_n .
\]
The graph closure of this core Hamiltonian is exactly the graph of
Theorem~\ref{thm:diag-energy} with $\lambda_n=n+1/2$, by joint
truncation. Energy measurement in the number basis gives outcome $n+1/2$
with probability $|x_n|^2$; omitted mass equals $\|Q_Nx\|^2$.
On the core $U(t)^*aU(t)=e^{-it}a$ and
$U(t)^*a^\dagger U(t)=e^{it}a^\dagger$. Recombining,
$Q(t)=Q\cos t+P\sin t$ and $P(t)=P\cos t-Q\sin t$.
For unit core input $x$, take the centered vectors
$u=(Q-\langle x,Qx\rangle)x$, $v=(P-\langle x,Px\rangle)x$.
The inner-product difference
$\langle u,v\rangle-\langle v,u\rangle=i$ follows from the commutator.
Cauchy--Schwarz gives $\Delta Q\,\Delta P=\|u\|\|v\|\ge1/2$.
At $x=e_0$ the centered vectors are $e_1/\sqrt2$ and $ie_1/\sqrt2$.
Hence equality holds. The commutator identity has been asserted on
a specified common domain throughout.

\subsection{Perturbing an unbounded Hamiltonian}

The next construction gives a state-dependent spatial cutoff for a
Hamiltonian that need not be diagonal in the given basis. It can be used
when the resolvent of the interacting operator has not been computed.

\begin{theorem}[Finite-band perturbation]\label{thm:band-dynamics}
Let $V$ have uniformly named entries, $V_{jk}=\overline{V_{kj}}$,
$V_{jk}=0$ for $|j-k|>w$, and a supplied integer bound $B\ge\|V\|$
on finite vectors. On the graph domain of $D$, set $H=D+V$.
This graph is self-adjoint and generates a PR-presented unitary group
$W(t)$. If $v$ is supported below $J$, $|t|\le T$ for an integer $T$,
and $K\ge J+mw$ with $m\ge1$, then
\begin{equation}\label{eq:galerkin}
 \|W(t)v-e^{-itP_KHP_K}v\|
 \le 2\|v\|\,(BT)^m/m!\le2\|v\|\,2^{4BT-m}.
\end{equation}
\end{theorem}
\begin{proof}
A finite input produces finite output under $V$; its norm bound extends
$V$ and its adjoint to all names. The adjoint relation for $H$ is the
one for $D$ after moving the bounded $V$ to the other side. Closedness
follows by subtracting $Vx$ from the second graph component. Joint
truncations for $D$ approximate $H$ in graph norm as well.

Write $D_K=P_KDP_K$, $V_K=P_KVP_K$, $H_K=D_K+V_K$.
The finite matrix exponential, extended by the identity on the
complement, is evaluated by \eqref{eq:exp-pr} using
$L\ge T(\Lambda(K)+B)$. Conjugate transpose and the group law follow by
absolutely convergent binomial convolution for commuting powers.

Here is a finite-product proof of the cutoff estimate. For $n\ge1$
put $h=t/n$, $A_K=e^{-ihD_K}$, $C_K=-ihA_KV_K$, and $F_K=A_K+C_K$.
For any finite self-adjoint $L$,
$\|e^{-ihL}-I\|\le |h|\|L\|$ and
$\|e^{-ihL}-I+ihL\|\le h^2\|L\|^2/2$ by integrating its unitary
series once and twice. Expansion to first order consequently gives
\[
 \|e^{-ihH_K}-F_K\|\le h^2(\Lambda(K)+B)^2,\qquad
 \|F_K\|^2\le1+h^2 B^2 .
\]
For the second bound expand $\|v-ihV_Kv\|^2$; the mixed term vanishes.
Telescoping $n$ factors yields
\[
 \|e^{-itH_K}-F_K^n\|
 \le T^2(\Lambda(K)+B)^2 e^{T^2 B^2/(2n)}/n.
\]
Every limit in $n$ at a \emph{fixed} cutoff therefore has a PR rate.

Expand $F_K^n$ from the left, choosing $A_K$ or $C_K$ at each
factor, stopping a branch immediately on choosing its $m$th $C_K$.
The complete branches with $r<m$ kicks apply $V_K$ at most $r$ times
to $v$. Between kicks there are diagonal phases. These branches agree
for all $K\ge J+mw$, since each kick increases the support by at
most $w$. Each stopped branch has $m$ kicks, diagonal phases, and a
rightmost block of unexpanded $F_K$'s. Its norm on $v$ is at most
$(|h|B)^m e^{T^2B^2/(2n)}\|v\|$. There are $\binom nm$ such branches
(with zero branches for $m>n$). Since
$\binom nm/n^m\le1/m!$, subtracting the two cutoff expressions
costs at most $2\|v\|(BT)^m e^{T^2B^2/(2n)}/m!$.
For any \emph{two} specified cutoffs send $n$ to infinity using the preceding
telescoping estimate with the larger of the two energy bounds. This proves the
pairwise version of \eqref{eq:galerkin}. The expansion is a finite
subset enumeration with at most $2^n$ full words. One may replace
the finite matrices by their rational approximants throughout, budgeting
a local error $\eta$ per factor by $n\eta(M+1)^n$ for input norm one
and factor bounds $M$ when $\eta\le1$. These estimates give the
finite-precision induction directly.

Given arbitrary $x$ and error $\epsilon$, approximate by finite $v$
within $\epsilon/8$, take
$m\ge k+4BT+4+C$ with an integer $C\ge\|v\|$ for $\epsilon=2^{-k}$, and evaluate
the chosen matrix exponential at $v$ within $\epsilon/8$.
The Cauchy estimate constructs the strong limit, uniformly in bounded
time, and proves \eqref{eq:galerkin}. It preserves norms. Products pass
to the strong limit: insert the fixed vector $W(s)x$ between the
two approximants and use their contraction bounds. The finite
composition law gives $W(t)W(s)=W(t+s)$, and $W(-t)$ is its inverse.
Strong continuity follows from approximation by one finite evolution
uniformly on a time interval.

For a finite $v$, all sufficiently large cutoffs agree on $H v$ and
$H^2v$. Twice integrating the finite unitary evolution gives
\[
 \|(W(h)-I)v\|\le |h|\|Hv\|,\qquad
 \|W(h)v-v+ihHv\|\le |h|^2\|H^2v\|/2.
\]
The first inequality extends to every graph input by joint truncation.
For such an $x$ the derivative remainder is bounded by $|h|$ times $2\|H(x-P_Nx)\|+|h|\|H^2P_Nx\|/2$. Its first term is
bounded by $2(\|Q_NDx\|+B\|Q_Nx\|)$.
Conversely, if the trajectory at $x$ has derivative $z$, differentiate
$\langle e_j,W(h)x\rangle=\langle W(-h)e_j,x\rangle$ at zero.
This gives $(iz)_j=\lambda_j x_j+(Vx)_j$. Hence $iz-Vx$ is the
required name of $Dx$, identifying the entire generator graph.
The group law now gives the differential equation on this domain
at every time, including invariance of the domain.
\end{proof}

For instance $V=g(S+S^*)$ for the unilateral shift $Se_j=e_{j+1}$
and real coupling $g$ has $w=1$, $B\ge2|g|$. Taking
$\lambda_j=(j+1/2)^2$ gives an unbounded anharmonic level sequence
with adjacent-level mixing. The cutoff estimate is independent of the
magnitude of the largest retained energy; that magnitude affects the
cost of the finite exponential.

\section{A particle on the line}
\label{sec:particle}

Continuous spectrum can be treated using the same arithmetic of finite
approximants. Take Gaussian-rational step functions of compact support on
the real line; a basic interval is $[u,v)$ with rational endpoints.
Refining to a common finite partition makes
\[
 \langle f,g\rangle=\sum_I |I|\,\overline{f_I}g_I
\]
a rational complex computation. The Hilbert space used here, denoted
$L^2(\mathbb R)$, is the regular metric completion of these step functions.
Values of a general vector at a point have no role in the definition.
One can identify this presentation with the coordinate Hilbert space
explicitly. Replacing endpoints by dyadics within $\delta$ costs at most
$\sqrt{2\delta}\sum_I|f_I|$ by the triangle inequality. On each integer
unit interval take its normalized indicator and the Haar differences
$|J|^{-1/2}(1_{J_{\rm left}}-1_{J_{\rm right}})$ for dyadic subintervals
$J$ of that unit interval. Orthogonality follows by cancellation on the
smaller interval. On each finite dyadic partition the indicator of either
half is a linear combination of the parent indicator and its difference.
Successive subdivision therefore transforms every compact dyadic step
into a finite Haar expansion preserving its squared norm, with finite
inverse. Integer intervals and their dyadic nodes admit a PR enumeration
with PR inverse (encode the signed integer and the depth/offset by pairing).
The finite transforms and density estimates extend both isometries to
the regular completions.

\subsection{Momentum, translations and the free Hamiltonian}

Let $R$ be a positive integer and write $E_R$ for multiplication by
$1_{[-R,R)}$. It is an orthogonal projection of norm at most one.
Every named vector has a PR cutoff for $\|(I-E_R)x\|$ by taking the
radius of a sufficiently accurate step approximation.
For a real multiplier $b(p)$ with supplied PR bounds $M_R$ and
continuity moduli on $[-R,R]$, bounded multiplication by $E_Rb$
is computed on steps by subdivision. Local oscillation $\eta$ costs
at most $\eta\|f\|$ on a step input $f$; the bounded extension
costs at most $M_R$ times its input error. Define the graph
\[
 (x,y)\in H_b \quad\Longleftrightarrow\quad
 E_Ry=(E_Rb)x\ \text{ for every positive integer }R.
\]
Suppose also that $b$ has uniformly PR point approximants on each such
interval. This graph is closed by boundedness of the local multipliers.
Compact steps lie in the domain. The graph pairs
$(E_Rx,E_Ry)$ approximate any graph pair, with both tails supplied
by the names. To identify the adjoint, test its identity against
arbitrary compact steps supported on $[-R,R)$. The bounded real
multiplier there is self-adjoint by the finite integral and passage to
the limit. Density of those steps gives exactly $E_Ry=(E_Rb)x$.
Conversely the identity on bounded support extends by truncating both
graph components. Thus $H_b$ is self-adjoint.

Multiplication by $e^{-itb(p)}$ yields a strongly continuous unitary
group. On a fixed compact interval the modulus of this multiplier follows
from $|e^{is}-e^{it}|\le |s-t|$; the input support controls the tail.
The two head--tail estimates of Theorem~\ref{thm:diag-energy} hold
with $P_N$ replaced by $E_R$, $\Lambda$ by $M_R$, and $Q_N$ by $I-E_R$.
Conversely, projecting a putative derivative onto $[-R,R)$ identifies
it with $-i(E_Rb)x$ using the bounded multiplier Taylor estimate.
This proves the strong Schr\"odinger equation with precisely the graph
domain displayed above.

For the momentum observable choose $b(p)=p$, and for the free
Hamiltonian of mass one choose $b(p)=p^2/2$. The latter group is
\[
  (U_0(t)f)(p)=e^{-itp^2/2}f(p).
\]
Momentum in a rational interval $I$ is measured by $f\mapsto1_I f$.
In a unit state the probability is $\|1_I f\|^2$. The energy interval
$[a,b]$, for rational $0\le a\le b$, is obtained by retaining the
two intervals with positive absolute momentum between $\sqrt{2a}$
and $\sqrt{2b}$. Endpoint inclusions represent the same vector.
These projections are effective: approximate both endpoints by rational
numbers within $\delta$, then use the step-height bound times the square
root of the total displaced length. For general input first approximate
by a step and use contraction. This supplies the sharp Born probabilities
even at an energy threshold.

\begin{proposition}[Continuous energy spectrum, with witnesses]
For each named $E\ge0$ and positive dyadic $\epsilon$ there is a
PR-prepared unit vector $v$ with $\|(H_{p^2/2}-E)v\|\le\epsilon$.
Every named solution of $H_{p^2/2}x=Ex$ equals zero, and
$H_{p^2/2}-E$ has dense proper range with explicit witnesses.
For $E\le-\eta$ with a supplied rational $\eta>0$ the bounded inverse
of $H_{p^2/2}-E$ exists and has bound $1/\eta$.
\end{proposition}
\begin{proof}
Put $s=\sqrt{2E}$ and take an integer $A\ge 2+s$.
Choose positive dyadic $\delta\le\min(1/4,\epsilon/(4A))$ and rational
$q$ with $|q-s|\le\delta$. Use the normalized indicator of
$[q-\delta,q+\delta)$. On its support $|p-s|\le2\delta$ and
$|p+s|\le2s+2\delta$, so its residual norm is at most
$2\delta(s+\delta)\le\epsilon$.

For a supposed eigenvector, approximate $x$ by a step $f$ within
$\gamma>0$ with height bound $C$. Use the union $J$ of the two intervals
of radius $2\delta$ about $q$ and $-q$, with $|q-s|\le\delta$.
Off $J$, both $|p-s|$ and $|p+s|$ are at least $\delta$.
On every compact truncation outside $J$, multiplication by
$|p^2/2-E|$ bounds the norm below by $\delta^2/2$, hence
the eigenvector equation forces that truncation to vanish. Passing
to the known spatial tail gives $x=1_Jx$. Its norm is bounded
by $\gamma+C\sqrt{8\delta}$. Taking
$\delta\le\gamma^2/(8(C+1)^2)$ proves $x=0$.

For density of the range, apply the same cut to a compact step $f$.
The vector $f1_{J^c}$ differs from $f$ by at most $C\sqrt{8\delta}$,
and dividing it by $p^2/2-E$ on its support gives a named graph
preimage. Indeed the reciprocal there is bounded by $2/\delta^2$;
on the finitely many retained compact intervals its Lipschitz bound
is at most $4D/\delta^4$ for support in $[-D,D]$, so it can be
sampled directly. Define its value as zero on the discarded intervals.
For properness let $g=1_{[s,s+1)}$, a named vector by endpoint
approximation. If $(H_{p^2/2}-E)x=g$, projection onto $[s,s+d)$
for rational $0<d\le1$ gives
$\sqrt d\le d(s+1)\|x\|$. This contradicts the choice of
$d$ with $\sqrt d(s+1)\|x\|\le1/2$, available from rational
upper bounds. Thus the asserted continuous-spectrum behavior has
finite-observation witnesses for density and for the obstruction to inversion.
For negative separated $E$, multiplication by $1/(p^2/2-E)$
has the stated bound and locally computable moduli; the graph partner
is the input plus $E$ times the output. The approximate eigenvectors
also rule out a supplied bounded inverse at any $E\ge0$.
\end{proof}

\subsection{Fourier transport with an explicit norm estimate}

Use $p$ for momentum and $x$ for position. For a finite momentum step on a disjoint partition
$f=\sum_j c_j1_{[a_j,b_j)}$ put
\[
 ({\cal F}f)(x)=(2\pi)^{-1/2}\int e^{ixp}f(p)\,dp,\quad
 S=\sum_j|c_j|,\quad A=\int |f|,\quad D_f=\int |p f(p)|\,dp .
\]
All displayed constants can be replaced by rational upper bounds.
The scalar constant $\pi$ is computable, for example from
$4\sum_{j=0}^{n-1}(-1)^j/(2j+1)$ with error at most $4/(2n+1)$.
The integral defining $\mathcal F f$ uses finite intervals. For $|x|\le R$
one may integrate the exponential Taylor polynomial term by term, with
\eqref{eq:exp-pr} and $L$ bounding $R\max_j(|a_j|,|b_j|)$.
The transform has height at most $A/\sqrt{2\pi}$ and Lipschitz
bound $D_f/\sqrt{2\pi}$. For rational $x$ with $|x|>0$, integration on each interval gives
$|{\cal F}f(x)|\le2S/(\sqrt{2\pi}|x|)$. Consequently
\begin{equation}\label{eq:fouriertail}
 \int_{|x|>R}|{\cal F}f(x)|^2\,dx\le {4S^2\over \pi R}.
\end{equation}
For example take integer $R\ge1+16S^2/\epsilon^2$; truncation then
costs at most $\epsilon/2$ in norm. On $[-R,R)$ use mesh and sample
errors making the local uniform error at most
$\epsilon/(4(R+1))$. Thus $\mathcal F f$ is a named vector, including
when the support contains $x=0$. The conjugate-sign transform
$\mathcal G$ has the same estimates.

\begin{theorem}[Plancherel for the step completion]\label{thm:pr-fourier}
$\mathcal F$ and $\mathcal G$ extend to mutually inverse PR isometries of
the represented Hilbert space. The conjugated Hamiltonian
$\mathcal F H_{p^2/2}\mathcal G$ acts as $-\frac12 d^2/dx^2$ on compact
piecewise-polynomial functions with continuous function and first
derivative, using their piecewise second derivative as the graph value
with the indicated factor.
\end{theorem}
\begin{proof}
We give a norm proof to establish the stability needed when passing
from finite input to an arbitrary Hilbert name. For a rational
$\alpha>0$, direct integration of the two half-lines gives
\[
 {1\over 2\pi}\int_{\mathbb R}e^{-\alpha|x|} e^{ixu}\,dx
       =K_\alpha(u):={1\over\pi}{\alpha\over \alpha^2+u^2}.
\]
The discarded tails on the left have absolute integral at most
$2e^{-\alpha R}/\alpha$ before the prefactor, and
$e^{-\alpha R}\le1/(1+\alpha R)$ suffices for a PR cutoff.
Multiplying the two copies of the finite-interval Fourier integral,
inserting this formula, and substituting $u=p-q$ gives
\begin{equation}\label{eq:abel-fourier}
 \int e^{-\alpha|x|}|{\cal F}f(x)|^2\,dx
     =\int K_\alpha(u)\langle f,T_u f\rangle\,du,\qquad
 (T_u f)(p)=f(p-u).
\end{equation}
Here are finite-sum details for these changes of variables. On
$|x|\le R$ and $|p|,|q|\le D$ containing the momentum support, split
the rectangle in $p,q$ at all step endpoints. The Lipschitz bounds of
$e^{-\alpha|x|}e^{ix(p-q)}$ in the three coordinates are
$\alpha+2D,R,R$. With maximal mesh $\eta$, the integrated error of
one rectangular sum is bounded by
$2R A^2(\alpha+2D+2R)\eta/(2\pi)$. Exchanging finite sums is exact.
Removing the $x$ cutoff costs at most $A^2 e^{-\alpha R}/(\pi\alpha)$.

For each pair of step intervals $I,J$, the substitution reduces to
\[
 \int_I\int_J K_\alpha(p-q)\,dq\,dp
 =\int K_\alpha(u)\,|I\cap(J+u)|\,du.
\]
Replace the kernel on $[-2D,2D]$ by rational mesh samples, to uniform
error $\eta/(\pi\alpha^2)$ plus sampling error, since
$|K_\alpha'|\le1/(\pi\alpha^2)$.
For one mesh interval $L$, the equality is equality of areas of rational
polygons. The linear map $(p,q)\mapsto(p,p-q)$ has determinant $-1$
and preserves triangle areas by the rational determinant formula. The
vertical slice length after the change is
$1_L(u)|I\cap(J+u)|$, a piecewise linear function whose integral is the
polygon area by summing trapezoids. The kernel-approximation error on
either side after summation is bounded by $A^2$ times its uniform error.
This proves \eqref{eq:abel-fourier} at arbitrary precision.

The kernel has integral one. Substitute $v=\alpha/u$ on $u\ge\alpha$.
Evenness gives the total integral as
$4\pi^{-1}\int_0^1(1+s^2)^{-1}\,ds=1$: integration of the finite
geometric identity leaves error at most $1/(2n+1)$ and yields the
alternating series defining $\pi/4$. Its mass on $|u|>\delta$
is at most $2\alpha/(\pi\delta)$.
Translations preserve the norm by translating the interval partition.
On each interval, its symmetric difference from its translate has length
at most $2|u|$, hence
$\|T_u f-f\|\le S\sqrt{2|u|}$. This estimate also constructs $T_u$
for real named $u$ by rational translations. It follows directly that
the right side of \eqref{eq:abel-fourier} differs from $\|f\|^2$ by
at most
\[
 \|f\|S\sqrt{2\delta}+{4\|f\|^2\alpha\over\pi\delta}.
\]
The left side differs from $\|{\cal F}f\|^2$ by at most
$\alpha R\,\|{\cal F}f\|^2+4S^2/(\pi R)$, using
$1-e^{-v}\le v$ for $v\ge0$. A bound for $\|{\cal F}f\|^2$ before
isometry is known is $(A^2+4S^2)/\pi$, by splitting at $R=1$.
For a requested error $0<\epsilon\le1$, take rational
$C\ge1+A+S+\|f\|$. One valid schedule is
\[
 R=\lceil32 C^2/\epsilon\rceil+1,\qquad
 \delta=\epsilon^2/(128 C^4),\qquad
 \alpha={\epsilon\delta\over64(R+1)C^2}.
\]
Indeed $\|{\cal F}f\|^2\le 5C^2$, and the four error terms above
sum to less than $\epsilon$. Each choice needs only rational
arithmetic, with optional bounded dyadic rounding
(for a rational $a/b>0$ with $a,b$ positive integers, $2^{-b}\le a/b$).
This proves $\|{\cal F}f\|=\|f\|$. The identical proof gives isometry
of $\mathcal G$.

An isometry sends a regular step name to a regular Hilbert name by
evaluating a more accurate input step and rounding the output within
the remaining budget. For compact steps $f,g$, integration over the
finite product of their supports gives
$\langle{\cal F}f,g\rangle=\langle f,{\cal G}g\rangle$ with our
inner-product convention. Extension by continuity gives the adjoint
identity everywhere. Preservation of inner products follows by
polarization from the norm identity. It implies
$\mathcal G\mathcal F=I$ and $\mathcal F\mathcal G=I$ by testing against all
step vectors. Finally two integrations by parts give
$\mathcal G(f'')(p)=-p^2\mathcal G f(p)$ for the stated compact
piecewise-polynomial $f$ (with named coefficients and certified
breakpoints). Interior boundary terms cancel by continuity of $f,f'$.
On each bounded momentum interval the transforms of $f,f''$ agree
with their integral formulas, by compact step approximation in both
$L^2$ and $L^1$ on the original bounded support. This proves the graph
equation, including the second component as a named Hilbert vector.
\end{proof}

Position in an interval $I$ now has momentum-space projection
$\mathcal G1_I\mathcal F$ and probability $\|1_I{\cal F}v\|^2$ in unit
input $v$. For $M_b f(p)=e^{ibp}f(p)$, direct substitution on steps
and contraction on all names give
\[
 T_aM_b=e^{-iab}M_bT_a,\qquad
 {\cal F}T_af(x)=e^{iax}{\cal F}f(x),\qquad
 {\cal F}M_bf(x)=({\cal F}f)(x+b).
\]
All identities are equalities of Hilbert names. The translation and
multiplication groups are strongly continuous by the preceding step
estimates. They therefore give the exponentiated canonical commutation
relations and both sharp position and momentum interval experiments
on a genuine continuous-spectrum system. In physical units the energy
phase is $e^{-itp^2/(2m\hbar)}$, the Fourier kernel is
$(2\pi\hbar)^{-1/2}e^{ixp/\hbar}$, and one supplies positive lower
bounds for $m,\hbar$.

\subsection{A particle in a bounded potential}

Let $H_0=\mathcal F H_{p^2/2}\mathcal G$ on position space. Supply a
bounded self-adjoint operator $V$ by its action on steps and an integer
bound $B\ge\|V\|$. A real bounded position potential with local PR
continuity moduli and rational sampling is such an input, as is a
bounded square well with rational endpoints by cutting intervals
exactly. Put $H=H_0+V$ on the graph domain of $H_0$.
Subtracting $Vx$ from the second component proves closedness, and the
adjoint identity reduces to that of $H_0$ by bounded self-adjointness
of $V$. Graph approximants to a finite step $f$ are explicit:
take $\mathcal F E_R\mathcal G f$, choosing $R$ from
\eqref{eq:fouriertail} for $\mathcal G f$. Its graph partner is
$\mathcal F((p^2/2) E_R\mathcal G f)+V\mathcal F E_R\mathcal G f$.
This gives exactly the density service in Section~\ref{sec:stone}.
We now construct the remaining resolvent data for unitary propagation.

Fix a rational $\rho>0$, either sign $\sigma\in\{-1,1\}$, and an
integer $L\ge\max(1,\rho,2B)$. In momentum coordinates
$R_0=(H_0+\sigma iL)^{-1}$ is the multiplier $(p^2/2+\sigma iL)^{-1}$,
transported by $\mathcal F$. It has bound $1/L$, and its graph
partner at input $f$ is $f-\sigma iL R_0 f$. The identity
\[
 H+\sigma iL=(I+V R_0)(H_0+\sigma iL)
\]
on graph inputs gives a two-sided inverse
$R=R_0\sum_{j\ge0}(-V R_0)^j$. The norm of $V R_0$ is at most $1/2$,
so the tail from $j=m$ is at most $2^{1-m}$ before the leading $R_0$.
The geometric identity for finite sums and this bound verify both inverse
equations. Every output is in the graph since $R_0$ is applied last.
Symmetry gives the sharper estimate $\|R\|\le1/L$ by
$\|(H+\sigma iL)v\|^2=\|Hv\|^2+L^2\|v\|^2$. Thus, with
$C=\sigma i(\rho-L)R$,
\[
 (H+\sigma i\rho)^{-1}
     = R\sum_{j\ge0}(-C)^j,\qquad \|C\|\le 1-\rho/L.
\]
With $b=\lceil L/\rho\rceil$, $(1-\rho/L)^b\le1/2$
since $(1-a)^b(1+ba)\le1$ for $0\le a\le1$, by finite induction.
The companion \texttt{ResolventTail.lean} proves its cleared-denominator
form $u^k(u+v+kv)\le(u+v)^{k+1}$ in the object calculus.
The tail from $j=bm$ is at most $(L/\rho)2^{-m}$, again before the
leading factor. The graph partner and both inverse identities follow
as above, and the inverse has bound $1/\rho$ by symmetry.
Products of a variable number of approximate operators here use fixed
per-factor precision, chosen by the finite telescoping estimate
following \eqref{eq:galerkin}; iteration on finite step codes is
ordinary primitive recursion. This supplies all imaginary-axis
resolvents, in particular at every positive integer scale, for the
resolvent construction of the unitary group. The Hamiltonian acts on the
compact polynomial test functions above by
$-\tfrac12 f''+Vf$, with the full Sobolev-type domain specified by
the transported multiplier graph.
